% GBFS France Audit Catalogue --- Scientific Data Descriptor submission draft
% Adapted from the Patterns version (gold_standard_patterns.tex).
% Format: Background & Summary | Methods | Data Records | Technical Validation
%         | Usage Notes | Code Availability
% Scientific Data accepts LaTeX submissions in standard article class.

\documentclass[11pt,a4paper]{article}

\usepackage[T1]{fontenc}
\usepackage[utf8]{inputenc}
\usepackage{lmodern}
\usepackage[margin=2.5cm]{geometry}
\usepackage{microtype}
\usepackage{amssymb,amsmath,amstext}
\usepackage{booktabs}
\usepackage{graphicx}
\graphicspath{{figures/}}
\usepackage{algorithm2e}
\usepackage{xcolor}
\usepackage{url}
\usepackage[colorlinks=true,linkcolor=blue,citecolor=blue,urlcolor=blue]{hyperref}
\usepackage{enumitem}
\usepackage{caption}
\usepackage{titlesec}
\usepackage{authblk}

\titleformat{\section}{\large\bfseries\sffamily\color{black!85}}{\thesection}{0.8em}{}
\titleformat{\subsection}{\normalsize\bfseries\sffamily\color{black!75}}{\thesubsection}{0.8em}{}

\newcommand{\ie}{\emph{i.e.}}
\newcommand{\eg}{\emph{e.g.}}
\newcommand{\etc}{etc.}

\title{\bfseries\sffamily\Large GBFS France Audit Catalogue: a versioned, audited
reference dataset of 46{,}307 bike-sharing stations from 123 French operators}

\author[1]{Rohan Foss\'e\thanks{Corresponding author: \texttt{rfosse@cesi.fr}}}
\author[2]{Ga\"el Pallares}
\affil[1]{CESI Engineering School, Montpellier, France}
\affil[2]{CESI LINEACT (EA 7527), Montpellier, France}
\date{}

\begin{document}
\maketitle

\noindent\textbf{Abstract.}
The General Bikeshare Feed Specification (GBFS) is the open standard
used by French micromobility operators to publish station-level data
on the national portal \emph{transport.data.gouv.fr}. Although the
publication is regulated and the syntax is shared, the semantics of
GBFS feeds vary substantially across operators, which limits their
direct use in research. We release the GBFS France Audit Catalogue
dataset, a versioned and audited national catalogue of 46{,}307
bike-sharing stations spanning 123 operators and 97 cities, enriched
with five reference sources (INSEE Filosofi, BAAC, BD~TOPO,
OpenStreetMap, FUB Barometer). The dataset is derived from the raw
GBFS feeds by a documented six-step purging pipeline that addresses
five recurring anomaly classes (A1--A5), reclassifies 30.9\,\% of the
raw stations, and is shipped together with its audit report, JSON
Schema, Frictionless Data Package descriptor, Croissant manifest and
Docker image for bit-exact reproduction. Internal robustness checks,
a baseline comparison, a characterization of the free-floating
subset showing that the same GBFS \texttt{capacity} field is
populated according to six incompatible semantics across French
operators, and a cross-country pilot on three non-French dock-based
systems (Bicing Barcelona ES, Oslo Bysykkel NO, Bergen Bysykkel NO)
establish the technical validity of the release. The dataset is archived on Zenodo under concept DOI
\texttt{10.5281/zenodo.20125460} and licensed under ODbL v1.0.

\vspace{0.5em}
\noindent\textbf{Keywords:} bike sharing; GBFS; data quality; FAIR;
reproducibility; open mobility data; reference dataset; France.

% ===================================================================
\section{Background \& Summary}\label{secBackground}

Bike-sharing systems (BSS) are a structural component of European
urban mobility, and their operational data are increasingly published
in the General Bikeshare Feed Specification (GBFS), an open ontology
maintained by the Open Mobility Foundation \cite{Mobi2024,OMF2023}. In
France, the 2019 Mobility Orientation Law (LOM, article L.~1115-1)
makes the publication of these feeds on the national open-data portal
\emph{transport.data.gouv.fr} a regulatory obligation, so that 123
operators currently publish GBFS data continuously.

While syntactic conformance with the GBFS schema is now widespread,
the semantic content of the published feeds varies substantially. A
systematic audit of the 123 French systems identifies five recurring
anomaly classes whose presence is sufficient to bias any cross-system
quantitative analysis based on the raw feed. They are (A1) out-of-domain
inclusion of car-sharing services under the GBFS schema; (A2)
placeholder capacities reported as constant non-zero values across all
stations of a system; (A3) structural over-capacity on free-floating
fleets that publish virtual anchor points as
\texttt{station\_information}; (A4) geospatial errors including
transposed coordinates and topological outliers; and (A5) coverage
incompatible with urban analysis (overseas, macro-regional). On the
French corpus, these classes affect 22 of the 123 systems and 30.9\,\%
of the raw stations (95\,\% bootstrap CI [30.5, 31.3]).

We release the \emph{GBFS France Audit Catalogue} dataset, which
re-publishes the French national catalogue after applying a
documented six-step purging pipeline that addresses A1--A5. The
released artefact contains 46{,}307 stations across 97 cities, of
which 5{,}442 are fully audited dock-based stations and 39{,}235 are
free-floating anchors retained with an explicit type label rather than
silently discarded. Each station is enriched with twelve derived
variables obtained by spatial join or administrative aggregation with
five reference sources: INSEE Filosofi (income, Gini, poverty),
BAAC/Cerema (cycling crash density), BD~TOPO (cycling infrastructure),
OpenStreetMap (facilities) and the FUB Cycling Barometer (perception).

The dataset is positioned as a research commons rather than a
proprietary product. It is versioned (SemVer, current release 1.0.0),
deposited on Zenodo with a concept DOI that always resolves to the
latest release, accompanied by machine-readable metadata (JSON Schema,
Frictionless Data Package, Croissant manifest, DCAT-AP record), and
licensed under the Open Database Licence v1.0. The full audit pipeline
is released as open source on GitHub under the MIT licence, together
with idempotence unit tests, a fixed-seed Docker image for bit-exact
reproduction, and the audit report.

The release supports four classes of downstream research:
composite-index work \cite{Fosse2026imd}; demand modelling with
graph-based neural networks \cite{Lin2018BSSGraph} whose false-station
nodes are a known A3 confound; spatial-equity analysis combining
station access with INSEE Filosofi; and multimodal integration with
GTFS \cite{MobilityData2024Validator}. A transferability pilot on
Bicing Barcelona, Oslo Bysykkel and Bergen Bysykkel
(Section~\ref{secValidation}) shows that A1--A5 generalises to other
European dock-based systems without threshold recalibration.

% ===================================================================
\section{Methods}\label{secMethods}

\subsection{Data sources}
The audit is applied to the GBFS feeds inventoried on
\emph{transport.data.gouv.fr} and on the MobilityData catalogue as of
2026-04-01: 142 candidate feeds, of which 14 fail technical ingestion
(unreachable URL, malformed JSON, or no \texttt{station\_information}
endpoint) and 5 are identified by class A5 as out-of-perimeter,
leaving 123 systems entering the purging pipeline. Each feed is
ingested via a thread-pooled HTTP client (\texttt{utils/gbfs\_collector.py})
that performs GBFS auto-discovery and persists a typed
\texttt{station\_information} snapshot alongside the system identifier
and the timestamp of the request.

The five reference sources used for contextual enrichment are listed
in Table~\ref{tab:sources}. Each source is integrated either by
spatial joining (point-in-polygon and buffer queries) or by
administrative aggregation (at the IRIS or municipal level),
depending on its native granularity.

\begin{table}[ht]
\centering
\caption{Reference sources used for contextual enrichment.}
\label{tab:sources}
\small
\begin{tabular}{@{}llll@{}}
\toprule
\textbf{Source} & \textbf{Granularity} & \textbf{Variables} & \textbf{Operation}\\
\midrule
INSEE Filosofi & IRIS / municipality &
   Median income/CU, poverty rate, local Gini & Aggregation \\
BAAC (Cerema)  & geolocated accident &
   Severe-crash density (500\,m buffer) & Spatial join \\
BD~TOPO        & polylines &
   Cycle-infrastructure linear (300\,m buffer) & Spatial join \\
OpenStreetMap  & polygons \& points &
   Facility density (transit stops, shops, services) & Spatial join \\
FUB Barometer  & municipal score &
   Subjective cycling perception (2017--2021) & Aggregation \\
\bottomrule
\end{tabular}
\end{table}

\subsection{Anomaly taxonomy}
Five anomaly classes were identified through an inductive-deductive
workflow over five iterations (2025-04 to 2026-01) starting from the
four data-quality dimensions of ISO/IEC~25012 \cite{ISO25012} and
enriched by reverse engineering of free-floating fleets. The classes,
their signatures and their incidence on the French corpus are listed
in Table~\ref{tab:anomalies}.

\begin{table}[ht]
\centering
\caption{Five GBFS anomaly classes on the French corpus. Incidence is
reported at the system level except for A4 which is intrinsically
station-level.}
\label{tab:anomalies}
\small
\begin{tabular}{@{}llll@{}}
\toprule
\textbf{Class} & \textbf{Name} & \textbf{Signature} & \textbf{Incidence}\\
\midrule
A1 & Out-of-domain inclusion & Car-sharing advertised as BSS & 14 systems\\
A2 & Placeholder capacity & Constant non-zero $c_i$, \eg $c = 100$ & 3 systems\\
A3 & Structural over-capacity & Conditional averaging on free-floating & 8 systems\\
A4 & Geospatial error & Transposed coords; $>3\sigma$ outliers; out-of-FR & $\sim 3.8\,\%$ of stations\\
A5 & Out of perimeter & System area $>50{,}000$\,km$^2$; overseas & 5 systems\\
\bottomrule
\end{tabular}
\end{table}

Class A3 is the most consequential and the least readily detectable.
Free-floating fleets advertise virtual stations whose
\texttt{capacity} is computed by \emph{conditional averaging} over
non-empty anchor points:
\begin{equation}
\bar{c}_{\text{profile}}
= \frac{\sum_{i\,:\,c_i > 0} c_i}{\#\{i\,:\,c_i > 0\}}
\;\neq\;
\bar{c}_{\text{actual}}
= \frac{1}{N}\sum_{i=1}^{N} c_i,
\label{eq:a3-bias}
\end{equation}
where $c_i$ is the declared capacity at anchor $i$. The ratio
$\bar{c}_{\text{profile}} / \bar{c}_{\text{actual}}$ is approximately
unity on dock-based systems but grows to one or two orders of magnitude
on free-floating fleets. Detection therefore proceeds on the
per-system capacity distribution rather than on schema conformance.

\subsection{Six-step purging pipeline}
The transition from the raw catalogue to the certified release
proceeds through six sequential transformations
(Algorithm~\ref{alg:purge}), implemented in
\texttt{utils/data\_loader.py} and logged in the audit report. Early
steps remove out-of-domain or coarse-grained anomalies (A1, A2, A4
perimeter); later steps require system-level statistics
($\bar c$, centroids, $\sigma$) and operate on a pre-cleaned set.

\begin{algorithm}[ht]
\SetAlgoLined
\DontPrintSemicolon
\KwIn{National catalogue $\mathcal{C}$ (123 systems); thresholds
$\sigma_{\max} = 3$, $N_{\min} = 20$.}
\KwOut{Certified set $\mathcal{G}$, rejected set $\mathcal{R}$, audit report.}
$\mathcal{G} \gets \emptyset$; $\mathcal{R} \gets \emptyset$\;
\ForEach{system $s \in \mathcal{C}$}{
    Ingest GBFS feed $\to \mathcal{S}_s$\;
    \textbf{S1.} If $s$ is car-sharing, relabel \texttt{carsharing} and continue\;
    \textbf{S2.} If $\mathrm{Var}(c) = 0$ and $c > 0$, force \texttt{free\_floating}\;
    \textbf{S3.} Recompute $\bar c_{\text{actual}}$ on all stations and reassign topologically\;
    \textbf{S4.} Keep $\varphi_i \in [41^\circ,52^\circ]$, $\lambda_i \in [-6^\circ,10^\circ]$\;
    \textbf{S5.} Drop $d(\mathbf{x}_i, \boldsymbol{\mu}_s) > \sigma_{\max} \cdot \sigma_s$\;
    \textbf{S6.} If $|\mathcal{S}_s^{\text{dock}}| < N_{\min}$, label \texttt{too\_small}\;
    Update $\mathcal{G}, \mathcal{R}$ accordingly\;
}
\Return $(\mathcal{G}, \mathcal{R}, \text{audit report}, \text{metadata})$\;
\caption{GBFS-to-Gold-Standard purging pipeline.}
\label{alg:purge}
\end{algorithm}

Four design properties are enforced by unit tests executed in
continuous integration: \emph{idempotence} (two successive applications
produce the same result), \emph{reversibility} (each excluded station
is preserved in \texttt{rejected\_stations.parquet} with its exclusion
motive), \emph{logging} (every step records counts in, counts out and
transformation type), and \emph{explicit parameterisation} (all
thresholds are exposed as hyperparameters and documented in the
schema).

\subsection{Threshold calibration}
The topological threshold $\sigma_{\max} = 3$ is a compromise between
robustness to GPS noise and the preservation of legitimate peripheral
extensions (\eg an off-centre TGV station or a university campus).
The single-dimension sensitivity sweep reported in
Section~\ref{secValidation} shows that the national Top-10 ranking is
invariant for $\sigma_{\max} \in [2, 4]$. The robustness threshold
$N_{\min} = 20$ dock-based stations is calibrated on the literature
on meshed networks \cite{Saberi2018}, below which density and
connectivity metrics become noise-dominated. Hybridisation buffer
radii (300\,m for BD~TOPO, 500\,m for BAAC) follow standard French
mobility-research practice.

% ===================================================================
\section{Data Records}\label{secRecords}

\subsection{Repository structure}
The Audit Catalogue is deposited on Zenodo under concept DOI
\texttt{10.5281/zenodo.20125460}. Each tagged release receives its
own version DOI, while the concept DOI always resolves to the latest
release. The current release is v1.0.0, dated 2026-05.

The deposit contains five primary files, summarised in
Table~\ref{tab:files}. Each file is described by a Frictionless Data
Package descriptor and listed in the Croissant manifest, so that the
release can be ingested directly by Google Dataset Search, Hugging
Face Datasets and Kaggle.

\begin{table}[ht]
\centering
\caption{Files included in the Audit Catalogue v1.0.0 release on Zenodo.}
\label{tab:files}
\small
\begin{tabular}{@{}llrl@{}}
\toprule
\textbf{File} & \textbf{Format} & \textbf{Size} & \textbf{Content}\\
\midrule
\texttt{stations\_gold\_standard\_final.parquet} & Apache Parquet & 7.8\,MB & 46{,}307 certified stations, 30 columns \\
\texttt{rejected\_stations.parquet}              & Apache Parquet & 1.4\,MB & Excluded entries with motive \\
\texttt{audit\_report.pdf}                       & PDF            & 0.6\,MB & Per-system pipeline trace \\
\texttt{systems\_catalog.csv}                    & CSV            & 0.05\,MB & 142 candidate feeds inventoried \\
\texttt{gold\_standard.croissant.json}           & JSON-LD        & 0.02\,MB & Croissant manifest with DOI bindings \\
\bottomrule
\end{tabular}
\end{table}

\subsection{Schema of the main Parquet file}
The main file (\texttt{stations\_gold\_standard\_final.parquet})
exposes 30 columns, documented in Table~\ref{tab:schema}: identifiers
(5), spatial and administrative (5), station description (4), audit
outputs (2), and seven thematic blocks from contextual enrichment.
Every column carries a stable name, a declared dtype, a source pointer
and an empirical completeness rate. Columns marked with $\dag$ are
introduced by the audit pipeline; columns marked with $\ast$ are
derived from external sources.

\begin{table}[ht]
\centering
\caption{Released schema of \texttt{stations\_gold\_standard\_final.parquet}.
``Compl.'' is the empirical completeness rate on the certified subset.}
\label{tab:schema}
\small
\begin{tabular}{@{}lllr@{}}
\toprule
\textbf{Column} & \textbf{Type} & \textbf{Source / Description} & \textbf{Compl.}\\
\midrule
\multicolumn{4}{l}{\textit{Identifiers (5)}}\\
\quad\texttt{uid}                & string  & Audit Catalogue primary key                 & 100\,\%\\
\quad\texttt{station\_id}         & string  & GBFS native identifier                    & 100\,\%\\
\quad\texttt{system\_id}          & string  & Operator-system identifier                & 100\,\%\\
\quad\texttt{system\_name}        & string  & Operator-system label                     & 100\,\%\\
\quad\texttt{source}              & string  & Feed URL / catalogue source               & 100\,\%\\
\midrule
\multicolumn{4}{l}{\textit{Spatial and administrative (5)}}\\
\quad\texttt{lat, lon}             & float64 & Geofiltered WGS84 coordinates             & 100\,\%\\
\quad\texttt{city}                 & string  & Normalised city name                      & 100\,\%\\
\quad\texttt{commune\_name}        & string  & INSEE commune label                       & 99.1\,\%\\
\quad\texttt{code\_commune}        & string  & INSEE commune code                        & 99.1\,\%\\
\quad\texttt{region\_id}           & string  & Administrative region                     & 100\,\%\\
\midrule
\multicolumn{4}{l}{\textit{Station description (4)}}\\
\quad\texttt{station\_name}        & string  & GBFS station name                         & 99.8\,\%\\
\quad\texttt{address}              & string  & GBFS address                              & 41.2\,\%\\
\quad\texttt{capacity}             & int     & Declared capacity (may be a placeholder)   & 96.8\,\%\\
\quad\texttt{n\_stations\_system}  & int     & Total stations in parent system            & 100\,\%\\
\midrule
\multicolumn{4}{l}{\textit{Audit outputs (2)$^\dag$}}\\
\quad\texttt{station\_type}        & enum    & \{docked\_bike, free\_floating, carsharing\} & 100\,\%\\
\quad\texttt{fetched\_at}          & datetime& Timestamp of the audited snapshot          & 100\,\%\\
\midrule
\multicolumn{4}{l}{\textit{Topography (2)$^\ast$}}\\
\quad\texttt{elevation\_m}                  & float64 & BD ALTI elevation                & 99.4\,\%\\
\quad\texttt{topography\_roughness\_index}  & float64 & Local relief amplitude           & 98.1\,\%\\
\midrule
\multicolumn{4}{l}{\textit{Cycling infrastructure (2)$^\ast$}}\\
\quad\texttt{infra\_cyclable\_km}           & float64 & BD TOPO cycle-lane linear (300\,m) & 97.6\,\%\\
\quad\texttt{infra\_cyclable\_pct}          & float64 & Share of dedicated right-of-way  & 96.9\,\%\\
\midrule
\multicolumn{4}{l}{\textit{Safety (1)$^\ast$}}\\
\quad\texttt{baac\_accidents\_cyclistes}    & int     & Severe-crash count (500\,m, 5\,y)  & 99.8\,\%\\
\midrule
\multicolumn{4}{l}{\textit{Multimodal access (2)$^\ast$}}\\
\quad\texttt{gtfs\_heavy\_stops\_300m}      & int     & Heavy-transit stops within 300\,m   & 98.7\,\%\\
\quad\texttt{gtfs\_stops\_within\_300m\_pct}& float64 & Share of accessible heavy-transit  & 98.7\,\%\\
\midrule
\multicolumn{4}{l}{\textit{Socio-economy (5)$^\ast$}}\\
\quad\texttt{revenu\_median\_uc}            & float64 & INSEE Filosofi median income/CU     & 95.4\,\%\\
\quad\texttt{gini\_revenu}                  & float64 & Local Gini index                    & 93.0\,\%\\
\quad\texttt{revenu\_d1}                    & float64 & First-decile income                 & 95.2\,\%\\
\quad\texttt{ecart\_interquar}              & float64 & Interquartile income spread         & 95.2\,\%\\
\quad\texttt{part\_menages\_voit0}          & float64 & Share of car-less households        & 95.4\,\%\\
\midrule
\multicolumn{4}{l}{\textit{Modal share (1)$^\ast$}}\\
\quad\texttt{part\_velo\_travail}           & float64 & Share of commute by bike (INSEE)    & 95.1\,\%\\
\bottomrule
\end{tabular}
\end{table}

\subsection{Geographic distribution and station-type breakdown}
The certified set contains 5{,}442 dock-based stations (fully audited),
39{,}235 free-floating anchors (typed but not rigorously audited; see
Usage Notes) and 1{,}630 car-sharing entries (relabelled and excluded
from BSS use). The geographic distribution across administrative
regions is given in Figure~\ref{fig:region-stations}.
\^Ile-de-France concentrates V\'elib'; Nouvelle-Aquitaine aggregates
many systems with a smaller mass of dock-based stations after the A3
correction.

\begin{figure}[ht]
\centering
\includegraphics[width=0.7\linewidth]{fig03_region_stations.pdf}
\caption{Certified stations by administrative region. Annotations
report the number of GBFS systems per region.}
\label{fig:region-stations}
\end{figure}

\subsection{Machine-readable manifests}
Four machine-readable descriptors accompany the Parquet file and are
shipped both in the Zenodo deposit and in the GitHub repository:
\begin{itemize}[leftmargin=*,itemsep=0.1em]
\item \textbf{JSON Schema} (\texttt{schema/gold\_standard.schema.json})
extends the GBFS v3 schema with the audit columns and the enrichment
blocks.
\item \textbf{Frictionless Data Package descriptor}
(\texttt{datapackage.json}) is consumable by any
\texttt{frictionless}-aware tool \cite{FrictionlessData2023}.
\item \textbf{Croissant manifest}
(\texttt{gold\_standard.croissant.json}) makes the dataset
discoverable by Google Dataset Search, Hugging Face Datasets and
Kaggle \cite{Croissant2024}.
\item \textbf{DCAT-AP record} \cite{DCATAP2021} is registered with
European open-data portals.
\end{itemize}

% ===================================================================
\section{Technical Validation}\label{secValidation}

Technical validity of the release is established by five independent
checks: (i) the empirical completeness of each derived variable;
(ii) two internal robustness checks on the rule-based audit;
(iii) a single-dimension sensitivity sweep on the topological
threshold; (iv) a baseline comparison against three published audit
strategies; and (v) a cross-country pilot on three non-French dock-based
systems.

\subsection{Completeness}
Empirical completeness per column on the certified subset is reported
in Table~\ref{tab:schema}. With the exception of the operator-supplied
\texttt{address} field (41.2\,\%, retained verbatim from the raw
feed), every column exceeds the 95\,\% completeness threshold and
most exceed 99\,\%. Missing values on enrichment variables are
concentrated on stations in IRIS units for which INSEE does not
publish a Filosofi record (typically very small or recently
constituted communes).

\subsection{Internal robustness}
The rule-based audit was challenged by two surrogate procedures on a
stratified random sample of 500 stations covering all 123 systems.

\paragraph{Orthogonal-coverage check.}
Three rule-based classifiers were implemented, each restricted to a
single signal source: a \emph{name-based} classifier seeing only
\texttt{system\_name} and \texttt{station\_name}; a \emph{capacity-
statistics} classifier seeing only the \texttt{capacity} column; and a
\emph{geospatial} classifier seeing only $(\varphi,\lambda)$. The
pairwise Cohen's $\kappa$ between the three is near zero
($\kappa < 0.10$); Fleiss $\kappa = -0.07$. This is the expected
signature of complementary coverage: each classifier sees only one
slice of the taxonomy. The recall of their union against the
integrated pipeline reaches $97.2\,\%$ at a precision of $91.5\,\%$,
with only 7 stations out of 500 flagged by all three. A1--A5 cannot
be reduced to a single heuristic.

\paragraph{Full-task agreement check.}
Three independent decision-tree classifiers were implemented, each
seeing all signals but applying them in a different priority order
(name-first, capacity-first, geospatial-first). On the same 500-
station sample, the three classifiers reach Fleiss $\kappa = 0.80$,
with pairwise Cohen's $\kappa \in [0.70, 1.00]$. The 93/500 stations
of disagreement concentrate on the A1/A3 frontier (free-floating
fleets whose name does not match the floating-pattern dictionary).
The lower bound $\kappa = -0.07$ (orthogonal) and the upper bound
$\kappa = 0.80$ (full-task) bracket the inter-annotator agreement
that three independent annotators with full access to the payload
would plausibly achieve.

\subsection{Threshold sensitivity}
The topological threshold $\sigma_{\max}$ was swept on
$\{2.0, 2.5, 3.0, 3.5, 4.0\}$ on the certified corpus. Jaccard
similarities against the reference at $\sigma_{\max} = 3.0$ are
$0.66, 0.79, 1.00, 1.00, 1.00$ respectively; Kendall's $\tau$ on the
national ranking is $0.92, 0.95, 1.00, 1.00, 1.00$. The national
Top-10 is invariant across the entire grid, which is the criterion
that matters for any downstream public-policy reading.

\subsection{Baseline comparison}
Three published audit strategies are evaluated against the Gold
Standard on the same 123-system corpus
(Table~\ref{tab:baselines}). The \emph{Raw GBFS} baseline ingests
\texttt{station\_information} as is; the \emph{vehicle\_type filter}
removes systems whose name advertises a non-bike fleet; the
\emph{MobilityData audit} adapts the canonical GTFS-realtime
schema-level rules to GBFS.

\begin{table}[ht]
\centering
\caption{Comparison of the Audit Catalogue against three baseline audit
strategies on the same French GBFS corpus. Brackets are 95\,\%
confidence intervals (Wald on Dock, bootstrap on
$\rho_{\text{FUB}}$ over cities $n=2{,}000$, Wilson on A1/A3).}
\label{tab:baselines}
\small
\begin{tabular}{@{}lcccc@{}}
\toprule
\textbf{Pipeline} & \textbf{Dock} & \textbf{$\rho_{\text{FUB}}$} &
\textbf{A1 ex.} & \textbf{A3 fix.}\\
\midrule
Raw GBFS                 & 0.12 \footnotesize[.12, .12] & 0.38 \footnotesize[.01, .68] & 0\,\% & 0\,\%\\
\texttt{vehicle\_type}    & 0.15 \footnotesize[.15, .15] & 0.19 \footnotesize[$-$.20, .52] & 82\,\% \footnotesize[59, 94] & 0\,\%\\
MobilityData audit       & 0.15 \footnotesize[.15, .15] & 0.19 \footnotesize[$-$.20, .52] & 82\,\% \footnotesize[59, 94] & 0\,\%\\
\textbf{Audit Catalogue}   & \textbf{1.00} & 0.25 \footnotesize[$-$.13, .59] & \textbf{100\,\%} \footnotesize[82, 100] & \textbf{100\,\%} \footnotesize[91, 100]\\
\bottomrule
\end{tabular}
\end{table}

A3 reclassification recall is $0\,\%$ for every baseline (Wilson
upper bound $9\,\%$): structural over-capacity is invisible to
syntactic validators. The \emph{vehicle\_type filter} and
\emph{MobilityData audit} baselines produce identical numbers on
this corpus because schema-level checks catch nothing that the name
filter does not already exclude.

\subsection{Downstream sensitivity}
The audit corrections have measurable downstream effects on city-level
rankings. The Bordeaux agglomeration aggregates five operators of
which one is the free-floating \emph{Pony} fleet
(Table~\ref{tab:bordeaux}). \emph{Pony} alone publishes 2{,}996
\texttt{station\_information} entries that the audit reclassifies as
free-floating, recomputing their actual capacity to 0.03 bike per
entry. The 225 stations of the historical \emph{V\'elo par TBM}
dock-based system remain. At the agglomeration level the nominal
docking total drops from 35{,}952 to 690, and the national IMD rank
(a soft-mobility composite from the companion study
\cite{Fosse2026imd}) moves from 2 to 14. The twelve agglomerations
most exposed to A3 are reported in Figure~\ref{fig:bordeaux-before-after}.

\begin{table}[ht]
\centering
\caption{Bordeaux: downstream sensitivity to the A3 correction at
the subsystem and agglomeration levels.}
\label{tab:bordeaux}
\small
\begin{tabular}{@{}lrr@{}}
\toprule
\textbf{Indicator} & \textbf{Raw GBFS} & \textbf{Audit Catalogue}\\
\midrule
\textit{Pony Bordeaux subsystem} & & \\
\quad \texttt{station\_information} entries & 2{,}996 & 2{,}996 (relabelled) \\
\quad Declared as dock-based                 & 2{,}996 & 0 \\
\quad Recomputed actual capacity              & --      & 0.03 bike/entry \\
\midrule
\textit{Bordeaux agglomeration (5 systems)} & & \\
\quad Total raw entries                      & 9{,}921 & -- \\
\quad Certified dock-based stations           & --      & 225 \\
\quad Nominal docking total (dock-based only) & 35{,}952 & 690 \\
\quad IMD rank \cite{Fosse2026imd}            & 2       & 14 \\
\bottomrule
\end{tabular}
\end{table}

\begin{figure}[ht]
\centering
\includegraphics[width=0.9\linewidth]{fig04_bordeaux_before_after.pdf}
\caption{Twelve French urban areas with the largest absolute
reduction in dock-based station count between the raw GBFS feed and
the certified Audit Catalogue.}
\label{fig:bordeaux-before-after}
\end{figure}

\subsection{Cross-country pilot}
Two non-French dock-based systems were audited with the same rule set,
recalibrating only the national perimeter (latitude/longitude bounds);
all other thresholds were kept at their French values
(Table~\ref{tab:e5-pilot}).

\begin{table}[ht]
\centering
\caption{Cross-country pilot: applying the A1--A5 rule set to three
non-French dock-based systems with French thresholds and country-
appropriate perimeter only. \emph{Ratio} is
$\bar c_{\text{profile}} / \bar c_{\text{actual}}$ from
Eq.~\ref{eq:a3-bias}.}
\label{tab:e5-pilot}
\small
\begin{tabular}{@{}lcccc@{}}
\toprule
\textbf{System} & \textbf{Country} & \textbf{Raw stns.}
                & \textbf{Capacity ratio} & \textbf{Anomalies}\\
\midrule
Bicing Barcelona       & ES & 544 & 1.006 & none of A1--A5 \\
Oslo Bysykkel          & NO & 266 & 1.000 & none of A1--A5 \\
Bergen Bysykkel        & NO &  97 & 1.000 & none of A1--A5 \\
\midrule
French corpus (ref.)   & FR & $\sim$67{,}000 & --- & 22 systems flagged \\
\bottomrule
\end{tabular}
\end{table}

Both systems pass the audit with zero structural anomalies. The
capacity-profile ratio is essentially unity on both feeds, consistent
with clean dock-based operation. Vehicle-type metadata is populated
cleanly (Bicing: \{\texttt{bicycle}, \texttt{scooter\_standing}\};
Oslo: \{\texttt{bicycle}\}), so A1 is not triggered. Perimeter and
topological filters yield only 1 outlier in Bicing, 3 in Oslo and 1
in Bergen.
These two negative controls do not establish full European
transferability, but they document that A1--A5 does not over-fit the
French market.

\subsection{Documented errors traceable to unaudited GBFS usage}\label{secConsequences}
A natural question for a Data Descriptor is whether the documented
anomalies have measurable downstream consequences, or whether they
are a theoretical concern that no real user has yet encountered. The
following four observations are drawn from the same audit on the
French corpus and document concrete errors that propagate when GBFS
is ingested without the present audit.

\paragraph{Semantic fragmentation of the same field across operators.}
On the 39{,}235 free-floating entries of the French corpus, the
GBFS \texttt{station\_information.capacity} field is populated
according to at least six mutually incompatible semantics
(Table~\ref{tab:capacity-semantics}). A downstream consumer who
applies a single arithmetic to the field (\eg \texttt{df.capacity.sum()})
aggregates six different quantities into one number. This is not a
theoretical hazard: it affects 84.7\,\% of the certified rows.

\begin{table}[ht]
\centering
\caption{Six observed semantics of the GBFS
\texttt{station\_information.capacity} field on the French
free-floating corpus, by operator brand. Total = 39{,}235 stations.}
\label{tab:capacity-semantics}
\small
\begin{tabular}{@{}llrrl@{}}
\toprule
\textbf{Operator} & \textbf{Cities} & \textbf{Stations} & \textbf{Reported $\bar c$} & \textbf{Semantics}\\
\midrule
Dott         & 8  & 20{,}224 & NaN  & Capacity left unfilled \\
Bird         & 11 &  4{,}499 & NaN  & Capacity left unfilled \\
Pony (Nice)  & 1  &     412  & 100  & Constant placeholder (A2) \\
Pony (Paris) & 1  &  3{,}617 & 1.6  & Per-vehicle ratio \\
Pony (Bordeaux, Angers, \dots) & 12 &  7{,}436 & 2--15 & Conditional fleet profile (A3) \\
Voi          & 4  &  2{,}816 & 5--10 & Small fleet-profile estimator \\
\bottomrule
\end{tabular}
\end{table}

\paragraph{Same standard, correct usage exists.}
The semantic fragmentation is not inherent to GBFS but operator-
specific. Beryl London (UK, hybrid free-floating) publishes
\emph{zero} entries in \texttt{station\_information} and reports all
vehicles in \texttt{free\_bike\_status}, which is the canonical
GBFS v2/v3 pattern for free-floating fleets. Beryl therefore avoids
the A3 anti-pattern entirely. The presence of a working positive
control demonstrates that the documented errors are propagated by
the absence of an audit tier between operator publication and
academic use, not by the GBFS schema itself.

\paragraph{Existing validators do not catch the errors.}
The official MobilityData GBFS Validator
\cite{MobilityData2024Validator} returns no error on any of the 22
systems flagged by A1--A5, because the anomalies are at the
semantic rather than syntactic level. The schema-level baseline
reported in Table~\ref{tab:baselines} reaches $0\,\%$ recall on A3,
which makes the gap empirically measurable: a production-grade
schema audit will pass a free-floating fleet labelled as
\texttt{station\_information}, and only a semantically-aware step on
the capacity distribution recovers it.

\paragraph{Downstream rankings are sensitive.}
The audit corrections reorder city-level rankings used by public
policy. The Bordeaux agglomeration moves from rank 2 to rank 14
in the national IMD soft-mobility ranking \cite{Fosse2026imd} after
the A3 correction on the \emph{Pony Bordeaux} subsystem alone
(Table~\ref{tab:bordeaux}). The Spearman correlation between dock
count and the FUB Cycling Barometer drops from
$\rho = 0.38$ (raw) to $\rho = 0.25$ (audited)
(Table~\ref{tab:baselines}); the raw correlation over-fits
the FUB perceptual signal by counting virtual stations as if they
were physical infrastructure. Any analysis that takes raw counts
as input and produces a city-level conclusion inherits these biases.

% ===================================================================
\section{Usage Notes}\label{secUsage}

\subsection{Limitations}\label{secLimitations}
Four limitations should be considered by any user.

\textbf{L1 -- National coverage.} The release is limited to
metropolitan France. Extending the audit to overseas territories or
to other European jurisdictions requires revisiting the
macro-regional surface threshold (A5) and the geofilter bounds (A4).
The cross-country pilot of Section~\ref{secValidation} demonstrates
that the rest of the protocol transfers without recalibration.

\textbf{L2 -- Static versus dynamic.} The v1.0 release audits
\texttt{station\_information} only. Auditing the
\texttt{station\_status} (real-time availability) feed requires a
distinct protocol; preliminary results on a 2-day pilot
(43 systems, 4{,}635 stations, 3.0M snapshots) identify dynamic
anomalies on 12.9\,\% of monitored stations. A companion paper
\cite{Fosse2026flows} develops the dynamic protocol.

\textbf{L3 -- Free-floating asymmetry.} The dock-based subset
(5{,}442 stations) is fully audited; the free-floating subset
(39{,}235 stations) is correctly typed but \emph{not rigorously
audited}. Free-floating entries are preserved in the release with an
explicit \texttt{station\_type = free\_floating} label, so that
downstream consumers can either exclude them with a single column
filter or run their own free-floating-native audit. The appropriate
framing of the v1.0 release is therefore \emph{Audit Catalogue for
dock-based French BSS, with explicit free-floating annotation}.

\textbf{L4 -- Operator coverage.} Some operators publish partial GBFS
feeds or are unavailable during the audit window; the inventory
completeness is not guaranteed and is re-evaluated at each release.

\subsection{Reuse examples}
The following three patterns are documented as runnable notebooks in
the released repository.

\paragraph{Example 1 -- Loading and filtering.}
Reading the certified Parquet and filtering to the dock-based subset
of a given city:
\begin{quote}\small\ttfamily
import pandas as pd\\
gs = pd.read\_parquet("stations\_gold\_standard\_final.parquet")\\
paris = gs[(gs.city == "Paris") \&\\
\ \ (gs.station\_type == "docked\_bike")]\\
len(paris)\ \ \#\ returns\ 1{,}507
\end{quote}

\paragraph{Example 2 -- Spatial-equity analysis.}
\begin{quote}\small\ttfamily
by\_city = (\\
\ \ gs[gs.station\_type == "docked\_bike"]\\
\ \ .groupby("city").agg(\\
\ \ \ \ n=("uid", "size"),\\
\ \ \ \ income=("revenu\_median\_uc", "median"))\\
)\\
rho = by\_city.n.corr(by\_city.income, method="spearman")\\
\#\ rho\ $\approx\ +0.39$\ (95\%\ bootstrap\ CI)
\end{quote}

\paragraph{Example 3 -- Identifying mobility deserts.}
\begin{quote}\small\ttfamily
q1 = gs.revenu\_median\_uc.quantile(0.25)\\
deserts = gs[(gs.station\_type == "docked\_bike") \&\\
\ \ (gs.revenu\_median\_uc < q1) \&\\
\ \ (gs.gtfs\_heavy\_stops\_300m == 0)]
\end{quote}
This three-condition filter recovers 1{,}041 stations
(19.1\,\% of the dock-based corpus) with a Q1 income threshold of
18{,}860\,EUR per consumption unit. Figure~\ref{fig:mobility-deserts}
reports the top fifteen contributing agglomerations.

\begin{figure}[ht]
\centering
\includegraphics[width=0.7\linewidth]{fig06_mobility_deserts.pdf}
\caption{Top fifteen French urban areas by number of dock-based
mobility-desert stations recovered by the three-condition filter
of Example 3.}
\label{fig:mobility-deserts}
\end{figure}

\subsection{Engineering footprint}
The released Parquet weighs $7.8$\,MB on disk; it loads in $0.23$\,s
into a \texttt{pandas} DataFrame with a memory peak of $15.9$\,MB
($\sim 200{,}000$ rows/s on a commodity laptop). Typical research
queries are sub-25\,ms: filtering to dock-based stations $5.2$\,ms,
grouping by city $3.7$\,ms, a Paris-bounding-box query $6.0$\,ms, a
merge with FUB 2023 (62 cities) $21.7$\,ms.

% ===================================================================
\section*{Code Availability}\label{secCode}

The audit pipeline is available at
\url{https://github.com/rohanfosse/bikeshare-data-explorer}, licensed
under MIT, with a \texttt{CITATION.cff} file and a Docker image
(\texttt{ghcr.io/rohanfosse/gold-standard:1.0.0}) for bit-exact
reproduction. The pipeline is implemented in 1{,}319 lines of code
across three Python modules (\texttt{utils/data\_loader.py},
\texttt{utils/gbfs\_collector.py},
\texttt{scripts/collect\_status.py}) exposing 35 public functions and
one collector class. The Python environment is pinned in
\texttt{requirements.txt}; random seeds are fixed
(\texttt{numpy=42}, \texttt{python=42}); idempotence and FAIR-conformance
unit tests are executed in continuous integration; the audit report
is regenerated automatically on every release. Releases are mirrored
to Zenodo via the GitHub--Zenodo integration.

% ===================================================================
\section*{Data Availability}\label{secData}

The GBFS France Audit Catalogue dataset is archived on Zenodo under
concept DOI \texttt{10.5281/zenodo.20125460} (current release 1.0.0,
2026-05). Each tagged release receives its own version DOI. The
dataset is licensed under the Open Database Licence v1.0 (ODbL-1.0);
the source code is licensed under MIT. Raw GBFS snapshots used as
input are archived with SHA-256 hashes in the same Zenodo deposit.
A Croissant JSON-LD manifest, a Frictionless Data Package descriptor,
a DCAT-AP record and a JSON Schema are shipped alongside the Parquet
file.

% ===================================================================
\section*{Acknowledgements}
The authors thank the producers of the open data that made this
research possible: GBFS operators, MobilityData, OpenStreetMap, Cerema,
ONISR, INSEE, FUB and the French GTFS community.

\section*{Author Contributions}
\textbf{R.F.}: Conceptualisation, Methodology, Software, Data
curation, Formal analysis, Writing -- original draft.
\textbf{G.P.}: Conceptualisation, Methodology, Supervision,
Writing -- review \& editing.

\section*{Competing Interests}
The authors declare no competing interests.

\bibliographystyle{unsrt}
\bibliography{references}

\end{document}
